% page 329 of the facsimile (image p-328 of the scan)
% block 167.6 x 237.7 mm ; left margin 34.4 mm
\pgc{12.700mm}{0.000mm}{
\l{\cel{53}321}
\l{}
\l{kuliso. I. (\sou{tekn.}) Suportilo kun ranuro interna, alonge qua peco}
\l{\cel{3}moviva povas irar e retrovenar sen obstaklo. - II. (1) Ranuro}
\l{\cel{3}facita sur la +stejo di teatro, alonge qua onu movas la frami}
\l{\cel{3}di la dekoruri. (2) Parto di la teatro qua esas dop la +stejo.}
\l{\cel{3}- III. En borso financala : la loko ube operacas la kurtajisti}
\l{\cel{3}qui ne havas permiso oficala. - IV. Rifalduro en la supra}
\l{\cel{3}parto di kurteno, en qua onu igas pasar stango - en la bordo}
\l{\cel{3}di vesto, por igar pasar kordono qua posibligas lua min-longigo}
\l{\cel{3}per plisuri. - DEFR.}
\l{}
\l{\sou{kulminar}. (\sou{netrans.}) I. (\sou{astron.}) Atingar, pasante ye la meri-}
\l{\cel{3}diano, sua alteso maxima super horizonto. - II. (\sou{metaf}.)}
\l{\cel{3}Esar (sociale) tre altagrade. - DEFIRS.}
\l{}
\l{\sou{kulpar}. (\sou{netrans., pri, pro)}. Faliar pri agar sua devo, koncerne}
\l{\cel{3}od etiko o prudenteso e habileso. -EFIS.}
\l{}
\l{\sou{kultar}. (\sou{trans.}) I. Honorizar deo suprega. - II. (\sou{metaf.})}
\l{\cel{3}Adorar ulu, ulo. - III. Agar la ceremonii extera per qui}
\l{\cel{3}la homi honorizas deo. - DEFIRS.}
\l{}
\l{\sou{kultelo}. Instrumento mikra-dimensiona, facita por tranchar o}
\l{\cel{3}sekar, e qua konsistas ye lamo muntita an mancho. - FI.}
\l{}
\l{\sou{kultivar}. (\sou{trans.}) I. Submisar (tero) ad ula labori por igar}
\l{\cel{3}lu fertila. - II. Submisar (la planti) a sorgi destinata}
\l{\cel{3}a plubonigar la efiko da naturo. - III. Explotar (ula planti)}
\l{\cel{3}por rekoltar lia produkturo. - IV. (\sou{metafore}) Kultivar la}
\l{\cel{3}arti pacala, kultivar la cienco, spirito, e c. - DEFIRS.}
\l{}
\l{\sou{kulturo.}\cel{2}La rezultajo de la kultivo di la arti, di la cienci,}
\l{\cel{3}di literaturo, di etiko. - DEFIRS.}
\l{}
\l{\sou{kumino}.\cel{2}(\sou{bot.}) Planto umbelifera, di qua la grani esas aromoza}
\l{\cel{3}(li saporas quale anizo). - L.\cel{1}\sou{carum carvi}. EFI.}
\l{}
\l{\textquotedbl{}\sou{kumis\textquotedbl{}}.\cel{2}Drinkajo fermentacita quan la nomadi di la centro-}
\l{\cel{3}regioni di Azia e di la sudo-regioni di Rusia, fabrikas ek}
\l{\cel{3}la lakto di la kavalini. -}
\l{}
\l{\sou{kumpreso}. Peco de linjo, faldita plura-duople, quan onu aplikas}
\l{\cel{3}sur korpo-parto malada, por ibe mantenar la pansuro, recevar}
\l{\cel{3}la puso, e c. - DEFIRS.}
\l{}
\l{\sou{kumular}. (\sou{trans.}) I. (jurisprudenco). Asemblar en sua persono}
\l{\cel{3}(plura yuri, plura autoritati). - II. Exercar plura funcioni}
\l{\cel{3}nacionala salario-implikanta. - F.}
\l{}
\l{\sou{kun}.\cel{2}I. Akompanata da... - II. Karakterizata da...Juntita a...}
\l{}
\l{\sou{kuneiforma}. (\sou{arkeol.}) Qua havas la formo di konio, di flecho-}
\l{\cel{3}pinto (kom ex. la literi kuneiforma di la antiqua Asiriani,}
\l{\cel{3}Medi e Persi). - DEFIRS.}
}
